\section*{Решения}
\addcontentsline{toc}{section}{Решения}

\begin{enumerate}

    \item \textbf{Решение (Демо 2026, №4):}\\
    Дана матрица билинейной формы $B = \begin{pmatrix} -1 & 1 \\ -3 & 4 \end{pmatrix}$ в старом базисе $e = (e_1, e_2)$.
    По условию заданы координаты векторов старого и нового базисов в некотором стандартном базисе. Найдем матрицу перехода $C$ от $e$ к $\hat{e}$.
    По определению, $(\hat{e}_1, \hat{e}_2) = (e_1, e_2) C$. Составим матрицы из координатных столбцов:
    $$E = \begin{pmatrix} 3 & 2 \\ 1 & -1 \end{pmatrix}, \quad \hat{E} = \begin{pmatrix} 4 & 1 \\ 3 & 1 \end{pmatrix}$$
    Отсюда $C = E^{-1} \hat{E}$. Вычислим обратную матрицу $E^{-1}$:
    $$E^{-1} = \frac{1}{-3 - 2} \begin{pmatrix} -1 & -2 \\ -1 & 3 \end{pmatrix} = \begin{pmatrix} 1/5 & 2/5 \\ 1/5 & -3/5 \end{pmatrix}$$
    $$C = \begin{pmatrix} 1/5 & 2/5 \\ 1/5 & -3/5 \end{pmatrix} \begin{pmatrix} 4 & 1 \\ 3 & 1 \end{pmatrix} = \begin{pmatrix} 2 & 3/5 \\ -1 & -2/5 \end{pmatrix}$$
    Закон изменения матрицы билинейной формы: $\hat{B} = C^T B C$.
    $$C^T B = \begin{pmatrix} 2 & -1 \\ 3/5 & -2/5 \end{pmatrix} \begin{pmatrix} -1 & 1 \\ -3 & 4 \end{pmatrix} = \begin{pmatrix} 1 & -2 \\ 3/5 & -1 \end{pmatrix}$$
    $$\hat{B} = \begin{pmatrix} 1 & -2 \\ 3/5 & -1 \end{pmatrix} \begin{pmatrix} 2 & 3/5 \\ -1 & -2/5 \end{pmatrix} = \begin{pmatrix} 4 & 7/5 \\ 11/5 & 19/25 \end{pmatrix}$$
    \textbf{Ответ:} $\hat{B} = \begin{pmatrix} 4 & 7/5 \\ 11/5 & 19/25 \end{pmatrix}$.

    \item \textbf{Решение (Демо 2026, №5):}\\
    Матрица квадратичной формы обязана быть симметричной. Если исходная матрица задана как $B = \begin{pmatrix} -1 & 1 \\ -1 & 4 \end{pmatrix}$, выделим её симметрическую часть $A$:
    $$A = \frac{B + B^T}{2} = \begin{pmatrix} -1 & 0 \\ 0 & 4 \end{pmatrix}$$
    Матрица перехода $C$ совпадает с найденной в предыдущей задаче: $C = \begin{pmatrix} 2 & 3/5 \\ -1 & -2/5 \end{pmatrix}$.
    Новая матрица квадратичной формы: $\hat{A} = C^T A C$.
    $$C^T A = \begin{pmatrix} 2 & -1 \\ 3/5 & -2/5 \end{pmatrix} \begin{pmatrix} -1 & 0 \\ 0 & 4 \end{pmatrix} = \begin{pmatrix} -2 & -4 \\ -3/5 & -8/5 \end{pmatrix}$$
    $$\hat{A} = \begin{pmatrix} -2 & -4 \\ -3/5 & -8/5 \end{pmatrix} \begin{pmatrix} 2 & 3/5 \\ -1 & -2/5 \end{pmatrix} = \begin{pmatrix} 0 & 2/5 \\ 2/5 & 7/25 \end{pmatrix}$$
    \textbf{Ответ:} $\hat{A} = \begin{pmatrix} 0 & 2/5 \\ 2/5 & 7/25 \end{pmatrix}$.

    \item \textbf{Решение (Проскуряков, №1244):}\\
    Дана квадратичная форма $Q(x) = x_1^2 + 2x_1 x_2 + 2x_2^2 + 4x_2 x_3 + 5x_3^2$.
    Её матрица формируется путем деления коэффициентов при смешанных произведениях пополам:
    $$A = \begin{pmatrix} 1 & 1 & 0 \\ 1 & 2 & 2 \\ 0 & 2 & 5 \end{pmatrix}$$
    Соответствующая симметричная билинейная форма (по формуле $B(x,y) = x^T A y$):
    $$B(x,y) = x_1 y_1 + x_1 y_2 + x_2 y_1 + 2x_2 y_2 + 2x_2 y_3 + 2x_3 y_2 + 5x_3 y_3$$

    \item \textbf{Решение (Проскуряков, №1249):}\\
    Дана квадратичная форма $Q(x) = x_1^2 - 2x_1 x_2 + x_2^2$.
    Аналогично, делим коэффициент при $x_1 x_2$ пополам и записываем матрицу:
    $$A = \begin{pmatrix} 1 & -1 \\ -1 & 1 \end{pmatrix}$$

    \item \textbf{Решение (Демо 2026, №6):}\\
    Квадратичная форма: $Q = -x_1^2 + x_2^2 + 5x_3^2 + 2x_1 x_2 - 2x_1 x_3$.
    Группируем слагаемые с $x_1$ и дополняем до полного квадрата:
    $$Q = -(x_1^2 - 2x_1 x_2 + 2x_1 x_3) + x_2^2 + 5x_3^2 = -(x_1 - x_2 + x_3)^2 + x_2^2 - 2x_2 x_3 + x_3^2 + x_2^2 + 5x_3^2$$
    $$Q = -(x_1 - x_2 + x_3)^2 + 2x_2^2 - 2x_2 x_3 + 6x_3^2$$
    Выделяем полный квадрат по $x_2$:
    $$2x_2^2 - 2x_2 x_3 + 6x_3^2 = 2(x_2^2 - x_2 x_3 + 0.25x_3^2) - 0.5x_3^2 + 6x_3^2 = 2(x_2 - 0.5x_3)^2 + 5.5x_3^2$$
    Канонический вид: $-y_1^2 + 2y_2^2 + \frac{11}{2}y_3^2$, где $y_1 = x_1 - x_2 + x_3$, $y_2 = x_2 - 0.5x_3$, $y_3 = x_3$.
    Выразим старые переменные через новые (матрица перехода $S$, $x = Sy$):
    $x_3 = y_3$; \quad $x_2 = y_2 + 0.5y_3$; \quad $x_1 = y_1 + y_2 - 0.5y_3$.
    Нормальный вид: $-z_1^2 + z_2^2 + z_3^2$.
    \textbf{Ответ:} Ранг $r=3$, положительный индекс $p=2$, отрицательный $q=1$. Преобразование: $x_1 = z_1 + \frac{1}{\sqrt{2}}z_2 - \frac{1}{\sqrt{22}}z_3$, $x_2 = \frac{1}{\sqrt{2}}z_2 + \frac{1}{\sqrt{22}}z_3$, $x_3 = \sqrt{\frac{2}{11}}z_3$.

    \item \textbf{Решение (Проскуряков, №1571):}\\
    Квадратичная форма: $Q = x_1^2 + x_2^2 + 3x_3^2 + 4x_1 x_2 + 2x_1 x_3 + 2x_2 x_3$.
    Выделяем полный квадрат по $x_1$:
    $$Q = (x_1 + 2x_2 + x_3)^2 - 4x_2^2 - x_3^2 - 4x_2 x_3 + x_2^2 + 3x_3^2 + 2x_2 x_3 = (x_1 + 2x_2 + x_3)^2 - 3x_2^2 - 2x_2 x_3 + 2x_3^2$$
    Выделяем полный квадрат по $x_2$:
    $$-3(x_2^2 + \frac{2}{3}x_2 x_3) + 2x_3^2 = -3(x_2 + \frac{1}{3}x_3)^2 + \frac{1}{3}x_3^2 + 2x_3^2 = -3(x_2 + \frac{1}{3}x_3)^2 + \frac{7}{3}x_3^2$$
    \textbf{Ответ:} Канонический вид $y_1^2 - 3y_2^2 + \frac{7}{3}y_3^2$. Замена: $y_1 = x_1 + 2x_2 + x_3$, $y_2 = x_2 + \frac{1}{3}x_3$, $y_3 = x_3$.

    \item \textbf{Решение (Проскуряков, №1574):}\\
    Квадратичная форма: $Q = x_1 x_2 + x_1 x_3 + x_2 x_3$. Так как квадратов нет, сделаем замену: $x_1 = u + v$, $x_2 = u - v$, $x_3 = w$.
    $$Q = (u^2 - v^2) + w(u + v) + w(u - v) = u^2 - v^2 + 2uw$$
    Дополним до полного квадрата по $u$:
    $$Q = (u^2 + 2uw + w^2) - w^2 - v^2 = (u + w)^2 - v^2 - w^2$$
    Пусть $y_1 = u + w$, $y_2 = v$, $y_3 = w$.
    \textbf{Ответ:} Канонический вид $y_1^2 - y_2^2 - y_3^2$. 
    Обратная замена: $x_1 = y_1 + y_2 - y_3$, $x_2 = y_1 - y_2 - y_3$, $x_3 = y_3$.

    \item \textbf{Решение (Проскуряков, №1586):}\\
    Требуется найти нормальный вид для формы из задачи №1571.
    Её канонический вид: $Q = y_1^2 - 3y_2^2 + \frac{7}{3}y_3^2$.
    Сделаем нормирующую замену координат: $z_1 = y_1$, $z_2 = \sqrt{3} y_2$, $z_3 = \sqrt{\frac{7}{3}} y_3$.
    \textbf{Ответ:} Нормальный вид $z_1^2 - z_2^2 + z_3^2$.

    \item \textbf{Решение (Демо 2026, №7 и №18):}\\
    Матрица формы: $A = \begin{pmatrix} 1 & -3 & -1 \\ -3 & 1 & 1 \\ -1 & 1 & 5 \end{pmatrix}$.
    Характеристическое уравнение: $\det(A - \lambda E) = -\lambda^3 + 7\lambda^2 - 36 = 0$.
    Корни: $\lambda_1 = -2, \lambda_2 = 3, \lambda_3 = 6$. Следовательно, канонический вид: $-2y_1^2 + 3y_2^2 + 6y_3^2$.
    Собственные векторы (ищем из $(A-\lambda E)x = 0$ и нормируем):
    Для $\lambda_1 = -2$: $v_1 = (1, 1, 0)^T \implies e_1 = \frac{1}{\sqrt{2}}(1, 1, 0)^T$.
    Для $\lambda_2 = 3$: $v_2 = (1, -1, 1)^T \implies e_2 = \frac{1}{\sqrt{3}}(1, -1, 1)^T$.
    Для $\lambda_3 = 6$: $v_3 = (-1, 1, 2)^T \implies e_3 = \frac{1}{\sqrt{6}}(-1, 1, 2)^T$.
    \textbf{Ответ:} Канонический вид: $-2y_1^2 + 3y_2^2 + 6y_3^2$. Ранг $r=3$, индексы инерции $p=2, q=1$.
    Матрица линейного преобразования $C$ (по столбцам): $\begin{pmatrix} 1/\sqrt{2} & 1/\sqrt{3} & -1/\sqrt{6} \\ 1/\sqrt{2} & -1/\sqrt{3} & 1/\sqrt{6} \\ 0 & 1/\sqrt{3} & 2/\sqrt{6} \end{pmatrix}$.

    \item \textbf{Решение (Проскуряков, №1370):}\\
    Матрица формы: $A = \begin{pmatrix} 1 & 2 & 2 \\ 2 & 1 & 2 \\ 2 & 2 & 1 \end{pmatrix}$.
    Собственные значения матрицы, состоящей из $1$ на диагонали и $2$ вне её: $\lambda_1 = 5$, $\lambda_2 = -1$, $\lambda_3 = -1$.
    Поскольку канонический вид при ортогональном преобразовании состоит из собственных значений, получаем результат.
    \textbf{Ответ:} Канонический вид: $5y_1^2 - y_2^2 - y_3^2$.

    \item \textbf{Решение (Проскуряков, №1374):}\\
    Матрица формы: $A = \begin{pmatrix} 2 & -2 & 0 \\ -2 & 1 & -2 \\ 0 & -2 & 0 \end{pmatrix}$.
    Характеристическое уравнение: $\det(A - \lambda E) = -\lambda^3 + 3\lambda^2 + 6\lambda - 8 = 0$. Корни: $\lambda_1 = 1, \lambda_2 = 4, \lambda_3 = -2$.
    Собственные векторы:
    $\lambda_1=1 \implies v_1 = (2, 1, -2)^T \implies e_1 = (2/3, 1/3, -2/3)^T$.
    $\lambda_2=4 \implies v_2 = (2, -2, 1)^T \implies e_2 = (2/3, -2/3, 1/3)^T$.
    $\lambda_3=-2 \implies v_3 = (1, 2, 2)^T \implies e_3 = (1/3, 2/3, 2/3)^T$.
    \textbf{Ответ:} Канонический вид: $y_1^2 + 4y_2^2 - 2y_3^2$. 
    Ортогональное преобразование: $x_1 = \frac{2}{3}y_1 + \frac{2}{3}y_2 + \frac{1}{3}y_3$, $x_2 = \frac{1}{3}y_1 - \frac{2}{3}y_2 + \frac{2}{3}y_3$, $x_3 = -\frac{2}{3}y_1 + \frac{1}{3}y_2 + \frac{2}{3}y_3$.

    \item \textbf{Решение (Демо 2026, №8):}\\
    Матрица формы $A = \begin{pmatrix} \alpha-1 & \alpha-1 & -\alpha \\ \alpha-1 & 2\alpha & -\alpha \\ -\alpha & -\alpha & \alpha-2 \end{pmatrix}$.
    Главные миноры по критерию Сильвестра:
    $\Delta_1 = \alpha - 1$.
    $\Delta_2 = (\alpha-1)(2\alpha) - (\alpha-1)^2 = \alpha^2 - 1$.
    $\Delta_3 = \det(A) = (\alpha+1)(2 - 3\alpha)$.
    Положительная определенность ($\Delta_1>0, \Delta_2>0, \Delta_3>0$): система $\{\alpha>1, \alpha^2>1, 2-3\alpha>0\}$ не имеет решений.
    Отрицательная определенность ($\Delta_1<0, \Delta_2>0, \Delta_3<0$): система $\{\alpha<1, \alpha^2>1, 2-3\alpha>0 \implies \alpha<2/3\}$. Решением является $\alpha < -1$.
    \textbf{Ответ:} При $\alpha \in (-\infty, -1)$ форма отрицательно определена. Форма никогда не является строго положительно определенной.

    \item \textbf{Решение (Проскуряков, №1598):}\\
    Матрица формы: $A = \begin{pmatrix} 5 & 2 & -1 \\ 2 & 1 & -1 \\ -1 & -1 & \lambda \end{pmatrix}$.
    Критерий Сильвестра для положительной определенности:
    $\Delta_1 = 5 > 0$.
    $\Delta_2 = 5 - 4 = 1 > 0$.
    $\Delta_3 = 5(\lambda - 1) - 2(2\lambda - 1) - 1(-2 + 1) = 5\lambda - 5 - 4\lambda + 2 + 1 = \lambda - 2$.
    Требуется $\Delta_3 > 0 \implies \lambda - 2 > 0$.
    \textbf{Ответ:} $\lambda > 2$.

    \item \textbf{Решение (Проскуряков, №1600):}\\
    Матрица формы: $A = \begin{pmatrix} 1 & \lambda & 5 \\ \lambda & 4 & 3 \\ 5 & 3 & 1 \end{pmatrix}$.
    Для строгой положительной определенности необходимо, чтобы \textit{все} диагональные миноры были положительны. 
    Рассмотрим главный минор второго порядка, образованный 1-й и 3-й строками/столбцами (условие обобщенного Сильвестра): 
    $M = \begin{vmatrix} 1 & 5 \\ 5 & 1 \end{vmatrix} = 1 - 25 = -24$.
    Поскольку один из главных миноров меньше нуля (независимо от $\lambda$), матрица принципиально не может быть положительно определенной.
    \textbf{Ответ:} Таких значений $\lambda$ не существует ($\lambda \in \varnothing$).

\end{enumerate}
